\section{Cosmic Dignity}

This is a brief outline and summary of the first part of ``The Individual and the Becoming of the World''\footnote{\url{http://www.gornahoor.net/Individual/index.htm}}. Although its point can not be made by means of a logical argument, perhaps a logical argument can lead one to see the point.

\paragraph{Quest for certainty involves direct intuition}


The problem is stated in terms of the quest for certainty. Following the traditions of the rishis of the Vedic period, Evola claims that certainty is achieved via an ``intuitive'' knowing, that is, a direct perception into the nature of reality. Again, like the rishis, such knowledge has its source in consciousness itself, since it is only the contents of my consciousness that I can know directly.

\paragraph{Aware of objects as representations in consciousness}


Starting with that observation, I notice that what is called the ``real'' or ``objective'' world is what is presented -- or re-presented -- in my consciousness. Objects, then, are whatever, in my experience, seems to resist me. For example, I see my desk, I notice its colours, it is firm to my touch. Yet it resists my direct will, so I call it a ``desk''. However, there is no rationale to posit the existence of a ``real'' desk, in addition to the ``desk'' represented in my consciousness, so that the former is somehow the ``cause'' of the desk in my consciousness.

\paragraph{Self is not an object but rather the constant in every act of consciousness}


Besides the objects represented in my consciousness, I become aware of my self as subject. This self, however, is not just another one of those objects present in my consciousness; rather it is the constant presence in the flow of phenomena in my mind. It is not experienced, then, as an object, but as the subject.

\paragraph{The world is the result of my will.}


I further note that I can act upon certain objects in my consciousness. This power to act and control I call my ``will''.

\paragraph{My will can be spontaneous.}


If the world is my self and its representations in consciousness, then clearly it is ``I'' who creates it, that is, who creates its appearance in my consciousness. This is easier to see in a dream; in a dream there are things and other people. Nevertheless, clearly it is my own mind that is creating the appearance of those other things and persons in my dream consciousness. It is doing so, however, ``spontaneously'', without the conscious intervention of my ``I''. Hence, even in my waking state, I am creating the world arising in my consciousness, albeit spontaneously.

\paragraph{Or it can be free and create knowingly and purposefully.}


When the will is free, it can evaluate various possibilities and then actualize one of them -- that is, bring it into manifestation. That is the definition of a magical act, hence the term ``magical idealism'' to name Evola's philosophical system.

\paragraph{The essence of an idea is the same as the existence of a thing}


The difference between the idea, or essence, of a thing and the existence of a thing is a question of degree.

\paragraph{It is a matter of power to bring an idea into manifestation.}


The measure of power is the capability to bring an idea into manifestation, that is, to alter reality in conformance with one's will.

As Evola put it in ``Essays on Magical Idealism'': ``What distinguishes magical idealism is its essentially practical character: its fundamental requirement is not to substitute one intellectual conception of the world with another, but rather to create in the individual a new `dimension' and a new depth of life.''

The new dimension involves understanding the world in a radically new way. Rather than the self as a passive observer of the world which produces a mental copy of itself in consciousness, the self instead is recognized as the creator of the world as its conscious act. In effect, one's whole worldview is reversed. The world is created from the inside out, that is, from ideas and spiritual qualities which then are manifested. It then becomes a matter of power as to which or whose ideas take hold. The world is then the battleground of spiritual forces, not simply as abstract ideals, but as a power struggle to bring those ideals into manifestation.

This the problem of knowing, that is, of self-realisation or the ``magical problem'' is really an ethical problem. By participation in this struggle, the individual becomes ``conscious of his task and his cosmic dignity.'' This revaluation of values, that is, morality as the creative act and not as strictly inhibitive, is Evola's fundamental point:

\begin{quotex}
Today the meaning of the moral as cosmic worth has been lost, that up to now, it has been reduced to nothing more than the corroboration and canonisation of deficiency, weakness, and fear -- of the plague of beautiful sentiments, of noble virtues, of holy ideals -- that is, of the fundamentally \emph{immoral}.
\end{quotex}


So the issue really is how to bring about this moment of self-realisation. Since it is not simply a philosophical argument, Evola then turns, in the second part of the essay, to the question of initiation.

\flrightit{Posted on 2008-08-11 by Cologero}

